\documentclass[11pt,letterpaper]{article}

\usepackage[utf8]{inputenc}
\usepackage[spanish]{babel}
\usepackage[margin=2.5cm]{geometry}
\usepackage{enumitem}
\usepackage{titlesec}
\usepackage{xcolor}
\usepackage{listings}

\definecolor{unicaucaverde}{RGB}{0,90,60}
\titleformat{\section}{\normalfont\Large\bfseries\color{unicaucaverde}}{\thesection}{1em}{}

\lstdefinestyle{cstyle}{
  language=C,
  basicstyle=\ttfamily\small,
  keywordstyle=\color{unicaucaverde}\bfseries,
  commentstyle=\color{gray}\itshape,
  stringstyle=\color{orange!70!black},
  numbers=left,
  numberstyle=\tiny\color{gray},
  frame=single,
  breaklines=true,
  showstringspaces=false,
  tabsize=2
}
\lstset{style=cstyle}

\title{\textcolor{unicaucaverde}{Ejercicios de refuerzo --- Semana 1}\\[4pt]
  \large Fundamentos de C: variables, tipos, operadores y entrada/salida}
\author{Programación --- Universidad del Cauca}
\date{2026}

\begin{document}
\maketitle

Práctica adicional para consolidar RA1.1 y RA1.2, para trabajo autónomo antes de la Semana 2. Cada
ejercicio incluye el enunciado y una solución completa en C.

\section*{Ejercicio 1 --- Conversión de temperatura}
Escribe un programa que lea una temperatura en grados Celsius y la convierta a Fahrenheit, usando
la fórmula $F = C \times 9/5 + 32$.

\begin{lstlisting}
#include <stdio.h>

int main(void) {
    float celsius, fahrenheit;

    printf("Temperatura en Celsius: ");
    scanf("%f", &celsius);

    fahrenheit = celsius * 9.0f / 5.0f + 32.0f;
    printf("%.1f C equivalen a %.1f F\n", celsius, fahrenheit);

    return 0;
}
\end{lstlisting}

\section*{Ejercicio 2 --- Promedio de tres notas}
Escribe un programa que lea tres notas (de 0.0 a 5.0) y calcule e imprima su promedio.

\begin{lstlisting}
#include <stdio.h>

int main(void) {
    float n1, n2, n3, promedio;

    printf("Ingrese 3 notas: ");
    scanf("%f %f %f", &n1, &n2, &n3);

    promedio = (n1 + n2 + n3) / 3.0f;
    printf("Promedio = %.2f\n", promedio);

    return 0;
}
\end{lstlisting}

\section*{Ejercicio 3 --- Interés simple}
Escribe un programa que lea un capital, una tasa de interés anual (en \%) y un tiempo en años, y
calcule el interés simple con $I = C \times r \times t / 100$.

\begin{lstlisting}
#include <stdio.h>

int main(void) {
    float capital, tasa, tiempo, interes;

    printf("Capital, tasa (%%) y tiempo (anios): ");
    scanf("%f %f %f", &capital, &tasa, &tiempo);

    interes = capital * tasa * tiempo / 100.0f;
    printf("Interes generado = %.2f\n", interes);

    return 0;
}
\end{lstlisting}

\section*{Ejercicio 4 --- Intercambiar dos variables}
Escribe un programa que lea dos enteros \texttt{a} y \texttt{b}, e intercambie sus valores usando
una variable temporal, mostrando los valores antes y después.

\begin{lstlisting}
#include <stdio.h>

int main(void) {
    int a, b, temp;

    printf("Ingrese a y b: ");
    scanf("%d %d", &a, &b);
    printf("Antes:   a = %d, b = %d\n", a, b);

    temp = a;
    a = b;
    b = temp;

    printf("Despues: a = %d, b = %d\n", a, b);
    return 0;
}
\end{lstlisting}

\section*{Ejercicio 5 --- Perímetro y área de un rectángulo}
Escribe un programa que lea la base y la altura de un rectángulo y calcule su perímetro
($2(b+h)$) y su área ($b \times h$).

\begin{lstlisting}
#include <stdio.h>

int main(void) {
    float base, altura, perimetro, area;

    printf("Base y altura del rectangulo: ");
    scanf("%f %f", &base, &altura);

    perimetro = 2 * (base + altura);
    area = base * altura;

    printf("Perimetro = %.2f\n", perimetro);
    printf("Area = %.2f\n", area);

    return 0;
}
\end{lstlisting}

\end{document}
